\section{Sparse-Deviation Decoding}
\label{sec:method}

% 证据：实验 1、4a、5；rulegen/ruledec.py、evalrule.py；neural/features.py、gp.py、nnrule.py、es.py；交接说明

Sparse-deviation decoding (\SDD) has three parts.
A \emph{construction rule} decides every shift by default.
A small number of \emph{deviation slots} in the genotype say at which shifts, and how, to depart from the rule.
A \emph{repair} step handles shifts where no robot is available.
The rule supplies the bias; the slots supply the freedom; the number of slots sets the balance between them.
This section describes the decoder (Sections~\ref{sec:method:rule}--\ref{sec:method:decode}),
the score used to learn the rule (Section~\ref{sec:method:proxy}),
three ways to learn it (Section~\ref{sec:method:acquisition}),
and how the number of slots is chosen (Section~\ref{sec:method:k}).

\subsection{Construction Rules}
\label{sec:method:rule}

A rule $\rul$ maps the decoding state at shift $i$ to a preference list of actions:
\begin{equation}
  \rul(s_i) = \bigl(\act_i^{(1)}, \act_i^{(2)}, \dots\bigr).
\end{equation}
The first action $\act_i^{(1)}$ is the \emph{default action}, taken when the shift is not deviated.
The later \reuse actions are the alternatives that a deviation can choose from, in the rule's order of preference.

The rule sees a read-only view of the state:
the current time and the remaining horizon;
for every available robot, its duty count, duty hours, remaining budget, number of duties left before maintenance, ready time and first-maintenance time;
the parameters of every robot type and how many robots of each type can still be opened;
and the capacity estimate of \dec{3}.
It also sees a \emph{default type} $\tau_g$, read from one gene, which it may use when it opens a new robot.
The decoder completes the list: available robots that the rule leaves out are appended in fewest-duties order, and $\tau_g$ is appended to the \new actions.
Every action therefore stays reachable, whatever the rule returns.

The simplest rule, which we call the \emph{default rule} $\rul_0$, sends the available robot with the fewest duties and opens a robot of type $\tau_g$ when none is available.
It is the ranking of \dec{2} without the genes.

\subsection{Genotype}
\label{sec:method:genotype}

The genotype has three parts:
\begin{equation}
  \genome = \bigl(\underbrace{p_1, a_1, \dots, p_\K, a_\K}_{\text{deviation slots}},\ \underbrace{g}_{\text{default type}},\ \underbrace{y_1, \dots, y_{\Nmax}}_{\text{tool genes}}\bigr).
\end{equation}
The tool genes are the same as in the baselines.
The number of slots is
\begin{equation}
  \K = \max\bigl(4,\ \lceil \dens \Nmax \rceil\bigr),
\end{equation}
where $\dens$ is the \emph{deviation density}.

Slot $k$ is a pair $(p_k, a_k) \in [0,1)^2$.
It is \emph{active} if $a_k \ge 0.6$; otherwise the decoder ignores it.
An active slot deviates at shift $i_k = \lfloor p_k \Nmax \rfloor$, and its value $u_k = (a_k - 0.6)/0.4 \in [0,1)$ selects the deviation:
\begin{itemize}
  \item If $u_k < 0.85$, the shift is \emph{reassigned}.
  The decoder takes the alternative robots of $\rul(s_{i_k})$ in the rule's order and picks the one at position $\lfloor (u_k/0.85)\,J \rfloor$, where $J$ is their number.
  If there is no alternative, the slot has no effect.
  \item If $u_k \ge 0.85$, the shift \emph{opens} a new robot of type $\lfloor (u_k - 0.85)/0.15 \cdot n_R \rfloor$.
  As in \dec{3}, the slot has no effect if the existing robots can cover the remaining horizon.
\end{itemize}
If several active slots point to the same shift, the last one is used.

Two choices in this design matter.
First, reassignment follows the rule's order. A small $u_k$ picks the rule's second choice, so even a deviation is biased toward actions that the rule rates highly.
Second, a slot is active with probability $\qon = 0.4$ under a uniform random genotype, so a random genotype deviates at about $\qon \K$ shifts.
The density $\dens$ thus sets how far from the rule a random genotype lands.
It also sets how many genotypes decode to the rule's own roster: all genotypes whose slots are inactive, a fraction $(1-\qon)^\K$ of the space.

\subsection{Decoding}
\label{sec:method:decode}

Algorithm~\ref{alg:decode} gives the decoder.
It first collects the active slots into a map $\Delta$ from shifts to deviation values.
It then builds the roster shift by shift.
At a shift with a deviation, the deviation is applied if it is valid.
At any other shift, the default action is taken.
If no robot is available at a shift without a deviation, the decoder runs the backtracking repair of \dec{5} before it opens a new robot.

\begin{algorithm}[t]
  \caption{Sparse-deviation decoding}
  \label{alg:decode}
  \begin{algorithmic}[1]
    \Require genotype $\genome$, rule $\rul$, density $\dens$, instance
    \Ensure roster, or \textsc{infeasible}
    \State $\K \gets \max(4, \lceil \dens \Nmax \rceil)$;\ \ $\Delta \gets \emptyset$
    \For{$k = 1, \dots, \K$}
      \If{$a_k \ge 0.6$} \State $\Delta[\lfloor p_k \Nmax \rfloor] \gets (a_k - 0.6)/0.4$ \EndIf
    \EndFor
    \State $\tau_g \gets \lfloor g\, n_R \rfloor$;\ \ $s \gets$ empty state;\ \ $i \gets 0$
    \While{$E(s) < \horizon$}
      \State $\avail \gets$ available robots in $s$
      \If{$\avail = \emptyset$ \textbf{and} $i \notin \Delta$}
        \State repair: change one of the last $W$ shifts
      \EndIf
      \State $(\act^{(1)}, \act^{(2)}, \dots) \gets \rul(s)$
      \State $\act \gets \act^{(1)}$
      \If{$i \in \Delta$ \textbf{and} the deviation $\Delta[i]$ is valid in $s$}
        \State $\act \gets$ deviation action (reassign or open)
      \EndIf
      \State simulate shift $i$ with action $\act$ and tool gene $y_i$
      \If{a robot or tool limit is exceeded} \Return \textsc{infeasible} \EndIf
      \State $i \gets i + 1$
    \EndWhile
    \State \Return roster in $s$
  \end{algorithmic}
\end{algorithm}

The key difference from the baselines lies in what happens after a change.
In \dec{0}--\dec{5}, changing the decision at shift $i$ changes the state, and every later gene is then read in the new state, often with a new meaning.
In \SDD, most later shifts have no gene at all: the rule decides them in the new state.
Only the few active slots are read again.
A change to one slot adds, removes or moves one deviation, and the rest of the roster adapts as the rule dictates.

\subsection{Proxy Score for Rule Quality}
\label{sec:method:proxy}

Section~\ref{sec:repr} showed that what matters is the distribution of rosters that random genotypes produce.
We therefore score a rule by that distribution, not by the single roster it builds.

For every training instance $I$ we fix a set $G_I$ of 67 genotypes:
three with all slots inactive (default-type gene $0.1$, $0.5$ and $0.9$) and 64 uniform random ones.
The same set is used for every candidate rule, so that differences between rules are not masked by sampling noise.
We decode every genotype with $\rul$ and compute the hypervolume of the non-dominated feasible rosters:
\begin{equation}
  \proxy(\rul) = \frac{1}{|\mathcal{I}_{\text{train}}|} \sum_{I \in \mathcal{I}_{\text{train}}} \HV_I\bigl(\{\mathcal{D}_\rul(\genome) : \genome \in G_I\}\bigr).
\end{equation}
$\HV_I$ uses a fixed normalization per instance:
$\fres$ is shifted by $\LB$, and both objectives are scaled by the ranges that the default rule $\rul_0$ produces on $G_I$;
the reference point is $(1.1, 1.1)$.
Scores of different rules are therefore directly comparable.
During rule learning we use $\dens = 0.25$.

One evaluation of $\proxy$ costs $15 \times 67 = 1005$ decodings, far less than one MOEA run per instance.
A rule with a high score is one around which random deviations find good rosters.
A rule that builds one good roster but has a poor neighbourhood scores low.
Section~\ref{sec:exp-static} shows that $\proxy$ ranks rules in the same order as full MOEA runs on unseen instances.

\subsection{Learning the Rule}
\label{sec:method:acquisition}

Because $\proxy$ only needs the rule as a black box, any rule representation can be optimized with it.
We use three, which cover hand-crafted features with symbolic search, learned weights, and free-form code.

\textbf{Features.}
The GP and neural rules score each candidate action from 27 scale-free features:
6 global (\eg remaining fraction of the horizon, capacity ratio, number of available robots),
7 for the robot type (\eg duty length, overhead per duty, duties per maintenance cycle),
9 for a \reuse candidate (\eg duty count and hours relative to the fleet mean, duties left before maintenance, whether maintenance can still finish before $\horizon$),
and 5 for a \new candidate (\eg robots of this type still allowed, whether it is the gene's default type).
All candidates are sorted by score, which gives the preference list.

\textbf{Genetic programming (GP).}
A rule is a pair of trees: one scores \reuse candidates, the other scores \new candidates.
We use a standard tree-based GP hyper-heuristic setup:
functions $\{+, -, \times, \div_{\text{protected}}, \max, \min\}$,
ramped half-and-half initialization (depth 2--6, maximum depth 8),
tournament size 7, crossover 0.80, mutation 0.15, reproduction 0.05, two elites,
population 64 and 30 generations.
Each generation is scored on a subset of 19 genotypes per instance (3 without deviations and 16 random);
the final rule is chosen with all 67.

\textbf{Neural network with evolution strategies (ES).}
A shared multilayer perceptron maps the 27 features of a candidate to a score.
It has one hidden layer of 16 $\tanh$ units (465 weights).
We first fit it by imitation: it learns to reproduce the decisions of good \dec{5} rosters (87\,451 state--action pairs).
We then optimize $\proxy$ directly with OpenAI-ES:
population 32, noise $\sigma = 0.02$, Adam with learning rate 0.01, 60 generations, and 23 genotypes per instance per generation.
Imitation alone is not enough, because good rosters disagree on what to do in the same state.
Its role is to start ES where $\proxy$ has a signal:
from random weights the network can open so many robots that every decoded roster falls outside the normalization box, and then $\proxy = 0$ for all perturbations.

\textbf{Large language model (LLM).}
The rule is a Python function that implements the interface of Section~\ref{sec:method:rule}.
We evolve it in the style of EoH \cite{liu2024evolution}.
Five initial rules are written from five prompts that describe different ideas.
Each of 12 generations creates four offspring (two refinements, one new idea, one combination of two parents) and keeps the six best rules.
The LLM sees each rule's score.
Generated code runs in a sandbox without imports, loops of unbounded length or access to private fields.
A run uses about 55 LLM calls and 55 evaluations of $\proxy$.

\subsection{Choosing the Deviation Density}
\label{sec:method:k}

The density $\dens$ sets the trade-off between bias and freedom.
Section~\ref{sec:theory} shows that its best value depends on how far good solutions are from the rule and on the evaluation budget.
These quantities are unknown for a new instance, so we choose $\dens$ once on the training instances.
For each candidate value and each learned rule, we run NSGA-II five times on all 15 training instances
and select the value with the highest mean HV.
The decision rule was written and tested on dummy data before the runs.
\todo{Final candidate set and selected $\dens$ (T10 selected $\dens = 1.0$ on the training instances, with 0.75 vs.\ 1.0 not significant; decide which value the paper uses).}
